\clearpage
\begin{summarybox}
\textbf{一句话：} 平面内到两个不同定点的\textbf{正距离比}固定，非等距时得到圆周，等距时得到中垂线。

\textbf{速查：} 以下圆公式要求 $a>0,\lambda>0,\lambda\ne1$。
\[
 A(-a,0),\ B(a,0):\quad
 h=\frac{a(\lambda^2+1)}{\lambda^2-1},\quad
 R=\frac{2a\lambda}{|\lambda^2-1|},\quad
 (x-h)^2+y^2=R^2.
\]
\[
 x_D=\frac{a(\lambda-1)}{\lambda+1},\quad
 x_E=\frac{a(\lambda+1)}{\lambda-1},\quad
 C=\text{$DE$ 中点},\quad R=\frac{DE}{2}.
\]
\[
 C=\frac{A-\lambda^2B}{1-\lambda^2},\qquad
 R=\frac{\lambda AB}{|1-\lambda^2|}
 \quad\text{（圆心式按坐标分别计算）}.
\]
\textbf{选择路线：} 数值题先平方配方；已知内外分点可直接找直径；一般坐标直接列距离平方式。
几何直觉是“内外角平分线垂直”，完整轨迹要有反向回验。

\textbf{最后四查：}
\begin{itemize}[leftmargin=*]
\item $A\ne B$？$\lambda>0$？$\lambda=1$ 是否单独处理？
\item 比值顺序对吗？$AB=2a$ 而不是 $a$？
\item 半径为正吗？圆心在正确的一侧吗？
\item 是否回到原比值、只取圆周，并保留所有附加限制？
\end{itemize}


\textbf{为什么会想到内外角平分线（发现过程四连）？}
\[
 \boxed{PA:PB\ \longrightarrow\ AD:DB=AE:BE\
 \longrightarrow\ PD\perp PE\ \longrightarrow\
 \text{以 $DE$ 为直径的圆}}
\]
\textbf{口诀：见比例，找分点；要证圆，造直角。}
$D$ 为内分点、$E$ 为外分点，位置均由比值唯一确定；
$\lambda>1$ 时 $E$ 在 $B$ 侧，$0<\lambda<1$ 时 $E$ 在 $A$ 侧。
若忘记圆心和半径，直接求 $D,E$，再取中点和半差；可用
$\boxed{x_Dx_E=a^2}$ 或 $h^2-R^2=a^2$ 做快速检验。
角平分线给出几何必要性，完整轨迹仍要有反向验证。
\end{summarybox}
